% TOP_PROBLEM: {"id":30007586,"problem_number":"LOCAL-30007586","title":"Causal mental models","category_id":21,"category":{"id":21,"name":"economics","display_name":"Economics","description":"Open economic theory statements, published research agendas, and proposed scoped specifications; question provenance and historical priority are qualified per record.","slug":"economics","order_index":21,"created_at":"2026-10-08T00:00:00Z"},"status":"research_question","external_url":null,"legacy_ids":[],"published":false,"statement":"Identify causal mental models from the specified belief and choice interventions, distinguishing causal understanding from ex-post rationalization and incentives.","economics":{"original_id":75,"field":"Behavioral and experimental economics","kind":"identification","formalization_basis":"adapted_research_question","source_status":"research_agenda","priority_status":"earliest_found","priority_note":"The 2025 full review verifies this methodological gap. It cites much older critiques of verbal explanations; first historical priority is not claimed.","earliest_verified_reference":{"authors":["Ingar Haaland","Christopher Roth","Stefanie Stantcheva","Johannes Wohlfart"],"title":"Understanding Economic Behavior Using Open-Ended Survey Data","year":2025,"url":"https://www.econtribute.de/RePEc/ajk/ajkdps/ECONtribute_362_2025.pdf","doi":"10.1257/jel.20251780","locator":"ECONtribute Discussion Paper 362, §2 introduction, printed p. 4; §3.1.1 Ex-post rationalization, printed p. 14; §5 Incentives, printed p. 40 (PDF pp. 6, 16, 42).","evidence_summary":"The review discusses retrospective rationalization, the possibility that contemporaneous elicitation changes deliberation, and an explicitly open question about incentives that improve truthful and accurate open-ended responses."},"current_reference":{"authors":["Ingar Haaland","Christopher Roth","Stefanie Stantcheva","Johannes Wohlfart"],"title":"Understanding Economic Behavior Using Open-Ended Survey Data","year":2025,"url":"https://www.econtribute.de/RePEc/ajk/ajkdps/ECONtribute_362_2025.pdf","doi":"10.1257/jel.20251780","locator":"ECONtribute Discussion Paper 362, §2 introduction, printed p. 4; §3.1.1 Ex-post rationalization, printed p. 14; §5 Incentives, printed p. 40 (PDF pp. 6, 16, 42).","evidence_summary":"The review discusses retrospective rationalization, the possibility that contemporaneous elicitation changes deliberation, and an explicitly open question about incentives that improve truthful and accurate open-ended responses."},"formalization_note":"The source explicitly discusses rationalization, elicitation reactivity, and truthfulness. The latent-graph identification framework and validation estimands are an adaptation."},"statusQualification":"Published question or research agenda verified at the cited locator. The mathematical specification is adapted for this catalogue and includes additional assumptions; source publication does not establish first-ever priority or certify this exact model remains unsolved. The source explicitly discusses rationalization, elicitation reactivity, and truthfulness. The latent-graph identification framework and validation estimands are an adaptation.","statusReviewedAt":"2026-10-08"}
% FIRST_POSTED: 2026-10-08
% ENTRY_KIND: statement_only
% !TeX program = lualatex
\documentclass[11pt]{article}
\usepackage{fontspec}
\usepackage[a4paper,margin=25mm]{geometry}
\usepackage{amsmath,amssymb,mathtools}
\usepackage{xurl}
\usepackage[hidelinks]{hyperref}
\newcommand{\sref}[2]{\hyperlink{source:#1:#2}{[S#2]}}
\setlength{\emergencystretch}{3em}
\begin{document}
% problemId:30007586
\section*{Causal mental models}\label{30007586}
\subsection{Definitions and mathematical statement}
Let individual \(i\)'s latent reasoning process be \(M_i\), represented by a directed acyclic graph over specified economic variables and a set of perceived conditional distributions. A task presents information \(X\), and the individual chooses \(A_i\). An elicitation protocol \(E\) produces text \(T_i\), which a fixed coding rule \(h\) converts into reported model \(\widehat M_i=h(T_i)\). Crucially, both \(A_i(E)\) and \(T_i(E)\) may depend on elicitation timing, incentives, and prompts. Consider randomized counterfactual-information interventions \(Z\) chosen to distinguish models in a prespecified finite class \(\mathcal M\). Define model predictive validity and elicitation reactivity as
\[V=\Pr(\widehat M_i\text{ predicts held-out responses to }Z),\qquad \rho=\mathbb E[A_i(E_1)-A_i(E_0)].\]
Here \(E_0\) omits explanatory prompts and \(E_1\) elicits explanations before or during choice, with all other task features fixed. Identify conditions under which observed text and intervention responses distinguish genuine pre-choice reasoning from retrospective justification, while estimating or bounding \(\rho\). High agreement between coders identifies reproducible classification, not necessarily latent-model accuracy. Multiple latent processes may yield the same text and choices, so point identification requires exclusions, informative interventions, or independent process measurements. The empirical objective is to establish validation criteria and uncertainty bounds for recovered models without assuming that verbal reports transparently reveal the causal mechanisms governing behavior.

\par\textbf{Formulation and scope.} The source explicitly discusses rationalization, elicitation reactivity, and truthfulness. The latent-graph identification framework and validation estimands are an adaptation.

\par\textbf{Open-question provenance.} An explicit question or research agenda was verified at the source locator. This does not by itself certify that every later version remains unresolved \sref{30007586}{1}.

\par\textbf{Historical priority.} The 2025 full review verifies this methodological gap. It cites much older critiques of verbal explanations; first historical priority is not claimed.

\subsection{Short English statement}
Identify causal mental models from the specified belief and choice interventions, distinguishing causal understanding from ex-post rationalization and incentives.

\subsection{Sources}
\begin{itemize}\raggedright
\item[\textnormal{[S1]}]\hypertarget{source:30007586:1}{} Ingar Haaland, Christopher Roth, Stefanie Stantcheva, Johannes Wohlfart. 2025. Understanding Economic Behavior Using Open-Ended Survey Data. ECONtribute Discussion Paper 362, §2 introduction, printed p. 4; §3.1.1 Ex-post rationalization, printed p. 14; §5 Incentives, printed p. 40 (PDF pp. 6, 16, 42)..
\url{https://www.econtribute.de/RePEc/ajk/ajkdps/ECONtribute_362_2025.pdf}.
DOI: 10.1257/jel.20251780.
{\small\textit{Earliest verified question/agenda reference; historical priority qualified below. The review discusses retrospective rationalization, the possibility that contemporaneous elicitation changes deliberation, and an explicitly open question about incentives that improve truthful and accurate open-ended responses.}}
\end{itemize}
\end{document}
